% ============================================================================
% AREAS FOR IMPROVEMENT
% ============================================================================

\section{Areas for Improvement}
\label{sec:improvements}

This section identifies opportunities for improving the \ctx{} codebase based on the analysis findings.

\subsection{High Complexity Functions}

\subsubsection{Problem}

The \func{extract\_symbols} functions in each parser handle too many responsibilities, leading to high fan-out counts (42--48 calls each).

\subsubsection{Current State}

\begin{lstlisting}[style=rustcode, caption={Current monolithic extract\_symbols (simplified)}]
pub fn extract_symbols(tree: &Tree, source: &str) -> Vec<Symbol> {
    let mut symbols = Vec::new();
    let query = Query::new(language(), SYMBOLS_QUERY)?;
    let mut cursor = QueryCursor::new();
    
    for match_ in cursor.matches(&query, tree.root_node(), source.as_bytes()) {
        // Symbol metadata extraction
        let name = get_capture_text(&match_, "name", source);
        let kind = determine_kind(&match_);
        
        // Signature building (complex logic)
        let signature = build_signature(&match_, source);
        
        // Docstring extraction
        let docstring = extract_docstring(&match_, source);
        
        // Visibility determination
        let visibility = determine_visibility(&match_);
        
        // Parent tracking
        let parent_id = find_parent(&match_, &symbols);
        
        symbols.push(Symbol { name, kind, signature, docstring, visibility, parent_id });
    }
    symbols
}
\end{lstlisting}

\subsubsection{Recommended Refactoring}

Decompose into smaller, focused functions:

\begin{lstlisting}[style=rustcode, caption={Proposed decomposition}]
pub fn extract_symbols(tree: &Tree, source: &str) -> Vec<Symbol> {
    let matches = run_symbols_query(tree, source);
    
    matches.into_iter()
        .map(|m| build_symbol_from_match(m, source))
        .collect()
}

fn run_symbols_query(tree: &Tree, source: &str) -> Vec<QueryMatch> {
    let query = Query::new(language(), SYMBOLS_QUERY)?;
    QueryCursor::new()
        .matches(&query, tree.root_node(), source.as_bytes())
        .collect()
}

fn build_symbol_from_match(match_: QueryMatch, source: &str) -> Symbol {
    Symbol {
        name: extract_name(&match_, source),
        kind: determine_kind(&match_),
        signature: build_signature(&match_, source),
        brief: extract_docstring(&match_, source),
        visibility: determine_visibility(&match_),
        parent_id: None, // Resolve in post-processing
    }
}
\end{lstlisting}

\subsection{Duplicate Formatter Code}

\subsubsection{Problem}

The \code{Formatter} trait implementations share significant code, particularly in \func{stream\_start} and \func{wrap} methods (85--88\% similarity).

\subsubsection{Recommended Solution}

Use default trait methods:

\begin{lstlisting}[style=rustcode, caption={Formatter with default implementations}]
pub trait Formatter {
    // Required methods
    fn name(&self) -> &'static str;
    fn content_wrapper(&self) -> (&'static str, &'static str);
    fn file_header(&self, path: &str) -> String;
    fn file_footer(&self) -> String;
    
    // Default implementations
    fn start(&self, tree_block: Option<&str>) -> String {
        let (start, _) = self.content_wrapper();
        match tree_block {
            Some(tree) => format!("{}\n{}\n", start, tree),
            None => start.to_string(),
        }
    }
    
    fn end(&self) -> String {
        let (_, end) = self.content_wrapper();
        end.to_string()
    }
    
    fn format_file(&self, path: &str, content: &str) -> String {
        format!("{}{}{}", 
            self.file_header(path),
            content,
            self.file_footer()
        )
    }
}
\end{lstlisting}

\subsection{Error Handling Consistency}

\subsubsection{Problem}

Error handling is inconsistent across modules:
\begin{itemize}
    \item Some functions use \code{io::Error}
    \item Others use \code{Box<dyn Error>}
    \item Embeddings use custom \code{EmbeddingError}
    \item Database uses \code{rusqlite::Error}
\end{itemize}

\subsubsection{Recommended Solution}

Create a unified error type using \code{thiserror}:

\begin{lstlisting}[style=rustcode, caption={Unified error type}]
use thiserror::Error;

#[derive(Error, Debug)]
pub enum CtxError {
    #[error("IO error: {0}")]
    Io(#[from] std::io::Error),
    
    #[error("Database error: {0}")]
    Database(#[from] rusqlite::Error),
    
    #[error("Parser error: {0}")]
    Parser(String),
    
    #[error("Embedding error: {0}")]
    Embedding(#[from] EmbeddingError),
    
    #[error("Query error: {0}")]
    Query(String),
    
    #[error("Configuration error: {0}")]
    Config(String),
}

pub type Result<T> = std::result::Result<T, CtxError>;
\end{lstlisting}

\subsection{Parser Abstraction}

\subsubsection{Problem}

Each parser implements similar patterns independently, leading to:
\begin{itemize}
    \item Code duplication
    \item Inconsistent behavior
    \item Difficulty adding new languages
\end{itemize}

\subsubsection{Recommended Solution}

Create a \code{LanguageParser} trait:

\begin{lstlisting}[style=rustcode, caption={LanguageParser trait}]
pub trait LanguageParser: Send + Sync {
    /// The tree-sitter language for this parser
    fn ts_language(&self) -> tree_sitter::Language;
    
    /// Query for extracting symbols
    fn symbols_query(&self) -> &str;
    
    /// Query for extracting call edges
    fn calls_query(&self) -> &str;
    
    /// Map capture names to symbol kinds
    fn kind_mapping(&self) -> &[(& str, SymbolKind)];
    
    /// Extract symbols from parsed tree (can override for custom logic)
    fn extract_symbols(&self, tree: &Tree, source: &str) -> Vec<Symbol> {
        // Default implementation using symbols_query()
        default_extract_symbols(self, tree, source)
    }
    
    /// Extract edges from parsed tree (can override for custom logic)
    fn extract_edges(&self, tree: &Tree, source: &str) -> Vec<Edge> {
        // Default implementation using calls_query()
        default_extract_edges(self, tree, source)
    }
}

// Factory function
pub fn parser_for(language: Language) -> Option<Box<dyn LanguageParser>> {
    match language {
        Language::Rust => Some(Box::new(RustParser)),
        Language::TypeScript => Some(Box::new(TypeScriptParser)),
        Language::Python => Some(Box::new(PythonParser)),
        Language::Solidity => Some(Box::new(SolidityParser)),
        _ => None,
    }
}
\end{lstlisting}

\subsection{Configuration Management}

\subsubsection{Problem}

Configuration is scattered across:
\begin{itemize}
    \item Command-line arguments
    \item Hardcoded values in source
    \item Environment variables (only \code{OPENAI\_API\_KEY})
\end{itemize}

\subsubsection{Recommended Solution}

Implement a configuration file system:

\begin{lstlisting}[style=bashcode, caption={Example .ctx/config.toml}]
# .ctx/config.toml

[index]
watch_debounce_ms = 500
batch_size = 50
exclude_patterns = ["*.generated.*", "vendor/"]

[embeddings]
provider = "local"  # or "openai"
dimension = 384
batch_size = 100

[output]
default_format = "xml"
show_sizes = false
include_tree = true

[complexity]
threshold = 10
warn_on_high = true
\end{lstlisting}

Configuration loading:

\begin{lstlisting}[style=rustcode, caption={Configuration loading}]
#[derive(Deserialize, Default)]
pub struct Config {
    #[serde(default)]
    pub index: IndexConfig,
    #[serde(default)]
    pub embeddings: EmbeddingsConfig,
    #[serde(default)]
    pub output: OutputConfig,
}

impl Config {
    pub fn load() -> Result<Self> {
        let config_path = Path::new(".ctx/config.toml");
        if config_path.exists() {
            let content = fs::read_to_string(config_path)?;
            Ok(toml::from_str(&content)?)
        } else {
            Ok(Config::default())
        }
    }
}
\end{lstlisting}

\subsection{Missing JSON Formatter}

\subsubsection{Problem}

The \func{get\_formatter} function has an incomplete JSON format implementation:

\begin{lstlisting}[style=rustcode]
OutputFormat::Json => Box::new(PlainFormatter), // TODO: Implement JSON formatter
\end{lstlisting}

\subsubsection{Recommended Implementation}

\begin{lstlisting}[style=rustcode, caption={JSON formatter implementation}]
pub struct JsonFormatter;

impl Formatter for JsonFormatter {
    fn start(&self, tree_block: Option<&str>) -> String {
        let tree = tree_block.map(|t| json!(t)).unwrap_or(json!(null));
        format!(r#"{{"project_tree": {}, "files": ["#, tree)
    }
    
    fn format_file(&self, path: &str, content: &str) -> String {
        let file = json!({
            "path": path,
            "content": content
        });
        file.to_string()
    }
    
    fn end(&self) -> String {
        "]}".to_string()
    }
}
\end{lstlisting}

\subsection{Summary of Recommendations}

\begin{table}[h]
\centering
\begin{tabular}{llc}
\toprule
\textbf{Issue} & \textbf{Solution} & \textbf{Priority} \\
\midrule
High complexity functions & Decompose into smaller functions & High \\
Duplicate formatter code & Default trait methods & Medium \\
Inconsistent error handling & Unified \code{CtxError} type & Medium \\
Parser code duplication & \code{LanguageParser} trait & Medium \\
Scattered configuration & Config file support & Low \\
Missing JSON formatter & Implement \code{JsonFormatter} & Low \\
\bottomrule
\end{tabular}
\caption{Improvement recommendations summary}
\label{tab:recommendations}
\end{table}
